
\section{应用与进一步发展}\label{sec:1-5}
\zhindex{yingyong}{应用}

本书有两章主要讨论应用：第~\ref{ch:applications-symplectic-topology} 章证明格罗莫夫 1985 年开创性论文中最早陈述的许多结果，
第 11 章讨论\BookRecordIndex{T0229}量子上同调。两章都包含主要结果的详细证明，并广泛讨论相关问题；这些问题在
索引中分别列于各主题之下，也列于“评注”和“例”两个总标题下。最易入手的应用包括：
第~\ref{sec:9-3} 节的\BookRecordIndex{T0120}非挤压定理证明；第~\ref{sec:9-4}、\ref{sec:9-5} 节关于辛 $4$-流形及其辛同胚群的结果；以及
第~\ref{sec:9-7} 节用\BookRecordIndex{T0143}格罗莫夫--威滕不变量区分辛结构的结果。阅读这些节只需对\BookRecordIndex{T0143}格罗莫夫--威滕
不变量 $\operatorname{GW}_{A,k}^M$ 有基本了解，不过其中确实会用到辛几何的柔性方法。
特别地，证明使用的是上面定理 A、B、C 的陈述，而不是其证明中的方法和思想。

第~\ref{sec:9-2} 节关于拉格朗日嵌入障碍的内容也相当直接。它研究带拉格朗日边界值的圆盘上一个
受扰\BookRecordIndex{T0209}柯西--黎曼方程的性质，因而需要稍多一些分析准备。另一方面，本节的第一个定理是
格罗莫夫的著名结果：欧氏空间上存在非标准辛结构。这个证明是格罗莫夫给出的简单直接论证，
以欧氏空间中\BookRecordIndex{T0219}拉格朗日子流形的性质为基础。在这一情形下不存在 $J$-全纯球面，因而紧致性
的讨论得到简化。

第~\ref{ch:applications-symplectic-topology} 章的其他应用更多地依赖第 8 章。第 8 章是关于\BookRecordIndex{T0156}哈密顿纤维化的预备章节，为研究
\BookRecordIndex{T0209}柯西--黎曼方程的扰动引入一个几何框架。受扰方程的解可以解释为平凡丛
$S^2\times M\to S^2$ 的截面；相对于某个适当的\BookRecordIndex{T0088}殆复结构 $\widetilde J_H$，这些截面是
全纯的。（$\widetilde J_H$ 与乘积\BookRecordIndex{T0088}殆复结构之差对应于扰动。）对有些应用（如第~\ref{sec:9-1} 节）
平凡丛已经足够，但若考虑任意黎曼曲面上的一般\BookRecordIndex{T0200}局部哈密顿纤维化，分析并不会更难。正如
第 6.7 节研究图像时那样，此时需要在更一般的框架中重新证明基本的\BookRecordIndex{T0408}正则性和紧致性结果。
论证本质上相同，不过一如既往，若要得到尖锐结果，少数棘手之处仍须小心处理。由此得到
研究哈密顿辛同胚群基本群的一个有用工具，即 \BookRecordIndex{T0047}Seidel 表示。还可以在定义域 $(\Sigma,j)$
的复结构固定时，定义高亏格\BookRecordIndex{T0143}格罗莫夫--威滕不变量。

第~\ref{sec:9-1} 节介绍这些思想在哈密顿动力学中的重要应用。第一个主要结果是：\BookRecordIndex{T0058}半正\BookRecordIndex{T0358}辛流形上的
每个哈密顿流至少有一条 $1$-周期轨道。要完整证明它，确实需要第 8 章的大部分内容；即使
只考虑平凡丛，仍须理解如何度量能量以及怎样处理\BookRecordIndex{T0243}冒泡。本节第二组结果证明，\BookRecordIndex{T0355}辛非球面流形
上存在两条不同的 $1$-周期轨道。由于假设排除了\BookRecordIndex{T0243}冒泡，这些结果略容易一些。第~\ref{sec:9-6} 节让
读者初步接触哈密顿群上的 Hofer 几何，介绍 Polterovich 和 Seidel 利用 \BookRecordIndex{T0047}Seidel 表示得到的
一个简单推论。这个表示真正只能在\BookRecordIndex{T0229}量子上同调或弗洛尔上同调的框架中定义。因此，第 11、
12 章都将回到它。

\paragraph{\BookRecordIndex{T0229}量子上同调。}
\zhindex{liangzishangtongdiao}{量子上同调}
第 11 章讨论\BookRecordIndex{T0229}量子上同调，它完全独立于第~\ref{ch:applications-symplectic-topology} 章；只有讨论 \BookRecordIndex{T0047}Seidel 表示的一节用到第 8 章。
因此，读完第 7 章即可直接阅读第 11 章。

\BookRecordIndex{T0229}量子上同调定义背后的基本思想很容易理解。为简单起见，假设 $(M,\omega)$（在球面意义下）
单调，即辛类 $[\omega]$ 在 $H_2(M)$ 的球面同调类上的限制，是第一陈类
$c_1(TM)$ 的一个正倍数 $\tau c_1(TM)$。再以 $H^*(M)$ 表示 $M$ 的整系数上同调模挠元。
因此可把 $H^*(M)$ 看成 $H^*(M,\mathbb Z)$ 在 $H^*(M,\mathbb R)$ 中的像。忽略挠元的
好处是，例如群 $H^k(M)$ 可以与对偶群
$\Hom(H^{2n-k}(M),\mathbb Z)$ 等同；换言之，一个 $k$ 次上同调类 $a$ 可以通过给出它与
所有 $c\in H^{2n-k}(M)$ 的配对值
\[
  \langle a,c\rangle:=\int_M a\smile c
\]
来描述。

对 $a\in H^k(M)$ 和 $b\in H^\ell(M)$，定义\BookRecordIndex{T0228}量子杯积 $a*b$ 为形式和
\begin{equation}\label{eq:1-5-1}
  a*b=\sum_{A\in H_2(M)}(a*b)_A q^{c_1(A)},
\end{equation}
其中 $q$ 是一个次数为 $2$ 的辅助变量，而上同调类
$(a*b)_A\in H^{k+\ell-2c_1(A)}(M)$ 由\BookRecordIndex{T0143}格罗莫夫--威滕不变量定义：
\[
  \langle(a*b)_A,c\rangle
  =\operatorname{GW}_{A,3}^M(a,b,c),
  \qquad c\in H^*(M).
\]
注意，类 $a,b,c$ 必须满足维数条件
\begin{equation}\label{eq:1-5-2}
  2n+2c_1(A)=\deg(a)+\deg(b)+\deg(c),
\end{equation}
\BookRecordIndex{T0143}格罗莫夫--威滕不变量 $\operatorname{GW}_{A,3}^M$ 才可能非零。因此 $c_1(A)\le2n$；又
由单调性 $0\le c_1(A)$，所以公式 \eqref{eq:1-5-1} 中只出现有限多个 $q$ 的幂。此外，
对 $q^d$ 的系数有贡献的类 $A$ 满足
$\omega(A)=\tau c_1(A)=\tau d$；由命题 4.1.5，其中只有有限多个能由 \BookRecordIndex{T0003}$J$-全纯曲线表示。
这说明 \eqref{eq:1-5-1} 右边的和有限。又因该和中只出现 $q$ 的非负整数次幂，故
$a*b$ 是群
\[
  \operatorname{QH}^*(M;\mathbb Z[q])=H^*(M)\otimes\mathbb Z[q]
\]
的元素，这里 $\mathbb Z[q]$
是以次数为 $2$ 的变量 $q$ 为变元的多项式环。线性延拓后，得到乘法
\[
  \operatorname{QH}^*(M;\mathbb Z[q])\otimes
  \operatorname{QH}^*(M;\mathbb Z[q])
  \longrightarrow\operatorname{QH}^*(M;\mathbb Z[q]).
\]
事实证明，这个乘法对加法具有分配律、分次交换且满足结合律。前两个性质显然，最后一个
性质则微妙得多，它依赖（分裂）公理。若 $(M,\omega)$ 不单调，\eqref{eq:1-5-1} 中的和
可能为无限和。为了赋予它意义，必须选择适当的系数环 $\Lambda$。\BookRecordIndex{T0231}量子系数环有若干种可能
选择；见例 11.1.4。

若 $A=0$，类 $A$ 中所有 \BookRecordIndex{T0003}$J$-全纯曲线都是常值曲线。因此
$\operatorname{GW}_{0,3}^M(a,b,c)$ 就是通常的三重交
$\int_M a\smile b\smile c$。对其他所有具有 $J$-全纯代表的 $A$，都有 $\omega(A)>0$，
所以 $a*b$ 的常数项恰为通常的杯积。换言之，$a*b$ 是通常杯积的一个形变。

例如，令 $M$ 为带有通常凯勒形式的复射影 $n$-空间 $\mathbb{CP}^n$。若 $p$ 是
$H^2(\mathbb{CP}^n)$ 的正生成元，$L\in H_2(\mathbb{CP}^n)$ 是由直线
$\mathbb{CP}^1$ 表示的类，那么任意两点间恰有一条直线这一事实反映在恒等式
\[
  \bigl\langle(p*p^n)_L,p^n\bigr\rangle
  =\operatorname{GW}_{L,3}^{\mathbb{CP}^n}(p,p^n,p^n)=1
\]
中。由 \eqref{eq:1-5-2}，其他类 $(p*p^n)_A$ 都因维数理由而消失。因此
$p*p^n=q\mathbf1$，其中 $\mathbf1\in H^0(\mathbb{CP}^n)$ 是单位元。还可证明，当
$i+j\le n$ 时 $p^i*p^j=p^{i+j}$。所以 $\mathbb{CP}^n$ 的\BookRecordIndex{T0229}量子上同调为
\[
  \operatorname{QH}^*(\mathbb{CP}^n;\mathbb Z[q])
  =\frac{\mathbb Z[p,q]}{\langle p^{n+1}=q\rangle}.
\]
例 11.1.12 给出直接证明。我们还说明怎样从 Batyrev 的环面簇\BookRecordIndex{T0229}量子上同调公式
（定理 11.3.4）\BookRecordIndex{T0326}推出这个结果；比较习题 11.3.11。这里字母 $p,q$ 的出现并非偶然。
第 11.3.2 节介绍 Givental 和 Kim 的一些近期工作：他们把\BookRecordIndex{T0274}旗流形的\BookRecordIndex{T0229}量子上同调环解释为
一个拉格朗日簇上的函数环。

第 11 章的目标之一，是展示\BookRecordIndex{T0229}量子上同调携带了多少信息。例如，对格拉斯曼流形和\BookRecordIndex{T0274}旗流形，
\BookRecordIndex{T0229}量子上同调反映 Verlinde 代数和 \BookRecordIndex{T0050}Toda 晶格等结构。它与其他数学领域的这些联系并不完全
出人意料，因为\BookRecordIndex{T0229}量子上同调构成镜像对称猜想的“A 侧”，而上述其他结构与“B 侧”相关。
我们不打算解释这句话的含义，但会讨论若干重要的相关概念，例如 \BookRecordIndex{T0020}Frobenius 流形和
\BookRecordIndex{T0015}Dubrovin 联络。还将说明怎样把\BookRecordIndex{T0143}格罗莫夫--威滕不变量中的信息编码进一个称为
\BookRecordIndex{T0144}格罗莫夫--威滕势的函数 $\Phi$。（分裂）公理等价于 $\Phi$ 满足一个称为 WDVV 方程的
三阶偏微分方程。

\paragraph{弗洛尔理论。}
最后一章介绍弗洛尔理论的基本定义和结构，但不深入分析细节。（其中许多内容已写入
Audin--Damian 的著作 \srcite{24}。）我们概述一个证明，说明弗洛尔理论中的环结构与
\BookRecordIndex{T0229}量子上同调中的环结构相同，并指出若干当前研究方向。特别地，我们定义 Schwarz 和 Oh
提出的哈密顿辛同胚\BookRecordIndex{T0271}谱不变量，并解释它们如何与 \BookRecordIndex{T0047}Seidel 表示相互作用。弗洛尔理论有许多
变体，其中一些在注 12.5.5 和第 12.6 节中概述。最后简要讨论\BookRecordIndex{T0356}辛涡旋方程；这些方程在
\BookRecordIndex{T0143}格罗莫夫--威滕不变量与规范理论之间架起桥梁。

\paragraph{结束语。}
本书讲解 \BookRecordIndex{T0003}$J$-全纯曲线理论的基础。辛拓扑还使用若干大不相同的技术，例如变分方法、
生成函数和 Donaldson 的殆复几何技术；本书讨论的椭圆型技术也衍生出许多理论，如弗洛尔
理论、辛同调和\BookRecordIndex{T0352}辛场论。一个自然的问题是：$J$-全纯球面理论处在这一谱系的什么位置？

纯粹的变分方法和生成函数方法，似乎只适用于底\BookRecordIndex{T0072}层空间类似余切丛、因而具有某种线性结构
的情况。例如，像 Hofer \srcite{177}、Viterbo \srcite{404} 或 Bialy--Polterovich
\srcite{37} 那样，可以用这些方法在欧氏空间的哈密顿辛同胚群
$\Ham(\mathbb R^{2n},\omega_0)$ 上定义不同版本的 \BookRecordIndex{T0028}Hofer 范数。Conley 和 Zehnder 在
\srcite{75} 中对环面的 Arnold 猜想所作的证明，先把问题提升到欧氏空间，再使用变分
方法。最近，Tamarkin \srcite{388} 以及 Guillermou--Kashiwara--Shapira \srcite{166}
使用\BookRecordIndex{T0073}层论方法（它是生成函数思想的重大推广）建立了一些基本的不可移位结果。这些方法也
适用于\BookRecordIndex{T0169}环面流形；如第 11.3 节所述，\BookRecordIndex{T0169}环面流形是欧氏空间一个开子集对一个复环面的商。
例如 Givental 在 \srcite{145} 中采用了这种途径。

处理一般\BookRecordIndex{T0358}辛流形时，看来必须使用 \BookRecordIndex{T0003}$J$-全纯曲线或某种版本的\BookRecordIndex{T0132}弗洛尔同调（可把后者看成
变分技术与椭圆型技术的结合）。例如，Lalonde--McDuff 在 \srcite{223} 中证明 \BookRecordIndex{T0028}Hofer
范数在任意\BookRecordIndex{T0358}辛流形上非退化时，使用了 $J$-全纯球面和辛嵌入技术。辛几何学家以许多巧妙
方式结合了现有方法。第 9、11、12 章会在适当位置提及其中一些可能性。
